\begin{proof}
Let 
\[
\mathbf a=(a_{1},\dots ,a_{n}),\qquad 
\mathbf b=(b_{1},\dots ,b_{n})\in\mathbb R^{n}.
\]
The Euclidean inner product and norm are  
\[
\langle\mathbf a,\mathbf b\rangle =\sum_{i=1}^{n}a_{i}b_{i},\qquad 
\|\mathbf a\|^{2}= \sum_{i=1}^{n}a_{i}^{2},\qquad 
\|\mathbf b\|^{2}= \sum_{i=1}^{n}b_{i}^{2}.
\]
We must prove \(\langle\mathbf a,\mathbf b\rangle^{2}\le\|\mathbf a\|^{2}\|\mathbf b\|^{2}\).

\medskip
\textbf{A non‑negative quadratic form.}  
For an arbitrary real parameter \(t\) consider  
\[
Q(t)=\sum_{i=1}^{n}(a_{i}-t\,b_{i})^{2}\ge 0\qquad(t\in\mathbb R).
\]

\medskip
\textbf{Expansion of \(Q(t)\).}  
Using \((a-tb)^{2}=a^{2}-2t a b+t^{2}b^{2}\),
\begin{align*}
Q(t)&=\sum_{i=1}^{n}\bigl(a_{i}^{2}-2t a_{i}b_{i}+t^{2}b_{i}^{2}\bigr)\\
    &=\Bigl(\sum_{i=1}^{n}b_{i}^{2}\Bigr)t^{2}
      -2\Bigl(\sum_{i=1}^{n}a_{i}b_{i}\Bigr)t
      +\Bigl(\sum_{i=1}^{n}a_{i}^{2}\Bigr).
\end{align*}
Set \(q=\sum b_{i}^{2}\), \(p=-2\sum a_{i}b_{i}\) and \(r=\sum a_{i}^{2}\).

\medskip
\textbf{Discriminant argument.}  
If \(q=0\) then \(\mathbf b=\mathbf 0\) and the inequality is trivial.  
Assume \(q>0\). Since \(Q(t)\ge 0\) for all \(t\), the discriminant of the quadratic
\(qt^{2}+pt+r\) must satisfy \(p^{2}-4qr\le 0\); that is,
\[
\bigl(-2\sum a_{i}b_{i}\bigr)^{2}
-4\Bigl(\sum b_{i}^{2}\Bigr)\Bigl(\sum a_{i}^{2}\Bigr)\le 0 .
\]

Dividing by \(4>0\) yields
\[
\bigl(\sum_{i=1}^{n}a_{i}b_{i}\bigr)^{2}
\le
\bigl(\sum_{i=1}^{n}a_{i}^{2}\bigr)
\bigl(\sum_{i=1}^{n}b_{i}^{2}\bigr),
\]
which is the Cauchy–Schwarz inequality.

\medskip
\textbf{Equality case.}  
Equality holds exactly when the discriminant is zero, i.e. when the quadratic
has a double root. Then \(Q(t)=q\,(t-t_{0})^{2}\) for some \(t_{0}\in\mathbb R\), and
\(Q(t_{0})=0\) forces \(a_{i}=t_{0}b_{i}\) for every \(i\). Hence there exists
\(\lambda=t_{0}\) such that
\[
a_{i}= \lambda\,b_{i}\qquad\text{for all }i,
\]
or one of the vectors is the zero vector. Thus equality occurs precisely when the
two vectors are linearly dependent.

\end{proof}